\begin{annexitem}{annex:p7.d.1}{thm:p07-t-simultaneous-recursion}{Arrow-only tuple recursion}
\begin{proof}
There is nothing to construct for an empty tuple. Otherwise set
\(\tau_i=0\to_{\mathrm{ty}}\sigma_1\to_{\mathrm{ty}}\cdots
\to_{\mathrm{ty}}\sigma_r\to_{\mathrm{ty}}\sigma_i\).
For each \(i\), use its own single recursor at type \(\tau_i\) to
form \(C_i(n)\) with equations
\[
\begin{split}
 C_i(0)&=\lambda k,\vec x.\,x_i,\\
 C_i(Sn)&=\lambda k,\vec x.\,
 C_i(n)(Sk,F_1(k,\vec x),\ldots,F_r(k,\vec x)).
\end{split}
\]
The second equation refers only to the previous \(C_i\), not to
any \(C_j\) being defined simultaneously. Its step term is
\(\lambda n,c,k,\vec x.\,c(Sk,\vec F(k,\vec x))\), so it is an
ordinary System-T definition with exactly the specified recursor type.
Put \(X_i(n)=C_i(n)(0,\vec a)\).

First define \(J(0)=\lambda k.k\) and
\(J(Sn)=\lambda k.J(n)(Sk)\) by a recursor at type
\(0\to_{\mathrm{ty}}0\). Open induction proves the functional
equality
\[
 \lambda k.J(n)(Sk)=\lambda k.S(J(n)(k)).
\]
Its zero case is beta conversion. At the successor, unfold \(J(Sn)\)
and apply the previous functional equality at \(Sk\); folding back
gives the displayed next instance. These are substitutions of whole
functionals, followed by definitional conversion. Evaluating at zero
shows \(J(Sn)(0)=S(J(n)(0))\), so open induction gives
\(J(n)(0)=n\).

The target proves, by open induction on \(n\), the following finite
conjunction of equalities at the types \(\tau_i\):
\[
 C_i(Sn)=\lambda k,\vec x.\,
 F_i(J(n)(k),C_1(n)(k,\vec x),\ldots,C_r(n)(k,\vec x)).
 \tag{F}
\]
At zero both sides reduce to \(\lambda k,\vec x.F_i(k,\vec x)\).
For the step, unfold \(C_i(S(Sn))\) and substitute (F) for the
whole functional \(C_i(Sn)\). Beta conversion gives
\[
 \lambda k,\vec x.\,
 F_i(J(n)(Sk),C_1(n)(Sk,\vec F(k,\vec x)),\ldots,
              C_r(n)(Sk,\vec F(k,\vec x))).
\]
The defining equations replace the latter arguments by
\(C_j(Sn)(k,\vec x)\); the defining equation
\(J(n)(Sk)=J(Sn)(k)\) gives the next instance of (F). The induction hypothesis is an equality
of functionals, so substitution inside this displayed term is legitimate.
No universal quantifier over \(k,\vec x\) has been inserted into an
open induction formula. Applying (F) to \(0,\vec a\) and using \(J(n)(0)=n\) gives the
required successor equations, and the zero equations are immediate.
All transformations are finite syntactic constructions.
\end{proof}
\end{annexitem}

\begin{annexitem}{annex:p7.d.2}{thm:p07-t-matrix-decision}{Numerical tests and contraction}
\begin{proof}
Define the numerical zero test by \(z(0)=0,z(Sn)=1\). Open induction
proves \(z(n)=0\leftrightarrow n=0\), and its values are zero and
one. At every finite type put
\[
 \operatorname{If}_\sigma(e,u,v)
   =R_\sigma\,u\,(\lambda k,r.v)\,e.
\]
The computation equations give its zero and successor values. With
\(p(0)=0,p(Sn)=n\), open induction gives
\(n=0\lor n=Sp(n)\); hence these cases cover every numerical input.
There is no equality test at type \(\sigma\).

Here is a numerical equality test with a direct open-induction proof,
so that no backward-challenge lemma is presupposed. For fixed \(m\),
use tuple recursion to define three numerical sequences
\((k_j,r_j,h_j)\), starting at \((0,m,0)\). If \(r_j=0\), let
\[
 (k_{j+1},r_{j+1},h_{j+1})=(k_j,0,Sh_j);
\]
otherwise let
\[
 (k_{j+1},r_{j+1},h_{j+1})=(Sk_j,p(r_j),h_j).
\]
The test in these definitions is just \(z(r_j)\). Open induction on
\(j\), with \(m\) unchanged, proves the quantifier-free invariant
\[
 j=k_j+h_j,\qquad m=k_j+r_j,\qquad r_j=0\lor h_j=0.
 \tag{D}
\]
At zero it is immediate. In the first branch the residual stays zero
and only \(h_j\) increases, giving the next first equation and
preserving the others. In the second branch
\(r_j=Sp(r_j)\) and the old disjunction forces \(h_j=0\).
Incrementing \(k_j\) and decrementing \(r_j\) preserves their sum
and increases \(k_j+h_j\) by one. These are the elementary addition
identities from its defining recursion.

Addition cancellation \(k+r=k+h\to r=h\) is proved by open
induction on \(k\), using \(Sk+r=S(k+r)\), itself proved from
addition's equations by induction on \(r\). Also
\(r+h=0\leftrightarrow(r=0\land h=0)\) follows by induction on
\(h\). Thus (D), at \(j=n\), gives
\[
 n=m\quad\longleftrightarrow\quad r_n=h_n
 \quad\longleftrightarrow\quad r_n=h_n=0.
\]
The second equivalence uses the last clause of (D). Consequently
\(E(n,m)=z(r_n+h_n)\) is the required equality test.
This construction uses only the previously proved tuple recursion,
zero tests, and ordinary open induction.

For Boolean formulas take the test of \(\neg A\) to be
\(\operatorname{If}(d_A,1,0)\), and the test of \(A\land B\)
to be \(\operatorname{If}(d_A,d_B,1)\); obtain the other
connectives from these. Structural induction proves zero exactly for
truth. Every atomic formula in an arithmetic Dialectica matrix is an
equality of numerical terms, even when those terms contain functional
parameters. This proves the required claim without higher-type
equality-testing constants.

In particular choose, component by component,
\[
 \vec h=\operatorname{If}(d_{A(\vec x,\vec y_1)},
                           \vec y_2,\vec y_1).
\]
If the first matrix is true, \(A(\vec x,\vec h)\) supplies the
second; if the first is false, \(A(\vec x,\vec h)\) itself is false.
Hence the target proves
\[
 A(\vec x,\vec h)\longrightarrow
 A(\vec x,\vec y_1)\land A(\vec x,\vec y_2).
\]
\end{proof}
\end{annexitem}

\begin{annexitem}{annex:p7.d.3}{thm:p07-t-backward-induction}{Changing the challenge}
\begin{proof}
Write \(p(0)=0\), \(p(Sn)=n\) for the predecessor recursor, and
\(a\mathbin{\dot{-}}0=a\),
\(a\mathbin{\dot{-}}Sb=p(a\mathbin{\dot{-}}b)\) for truncated
subtraction. Use \(a\leq b\) as the quantifier-free abbreviation
\(a\mathbin{\dot{-}}b=0\). The elementary numerical identities
needed below are
\[
\begin{gathered}
 a\mathbin{\dot{-}}a=0,\qquad
 Sa\mathbin{\dot{-}}Sb=a\mathbin{\dot{-}}b,\\
 Sw\leq u\ \Longrightarrow\
 w\leq u\ \land\
 u\mathbin{\dot{-}}w=S(u\mathbin{\dot{-}}Sw).
 \tag{N}
\end{gathered}
\]
They are open-arithmetic consequences of the displayed recursions and
successor axioms. The second identity follows by induction on \(b\),
and then the first by induction on \(a\). For the last assertion,
cancel a common successor in the two subtraction terms until one
argument is zero. Its zero cases are: \(Sw\leq0\) is false, and
\(S0\leq Su\) gives \(Su=S(u\mathbin{\dot{-}}0)\).
This cancellation argument is formalized by induction on a bound for
the sum of the two arguments. Its bounded list of cases can be encoded
as a finite Boolean conjunction, with conjunction computed by recursion;
open induction on that computation yields the same bounded induction.
The Boolean encoding is available from
\cref{thm:p07-t-matrix-decision}, whose direct countdown proof does not
use the present lemma. A finite conjunction is computed by
\(c(0)=0\) and \(c(Sj)=\operatorname{If}(c(j),d_j,1)\), where
\(d_j\) is the zero-for-truth test of its next case. The recursion
equations, with open induction, give extraction of each earlier case
from \(c(j)=0\) and the introduction of the next case. This supplies
the bounded case induction just used.

Fix the desired numerical parameter \(u\) and challenge \(\vec y\).
By \cref{thm:p07-t-simultaneous-recursion} define
\[
 \vec e(0)=\vec y,\qquad
 \vec e(Sz)=\vec b(u\mathbin{\dot{-}}Sz,\vec e(z)).
\]
We prove by open induction on \(w\) the single quantifier-free formula
\[
 w\leq u\ \longrightarrow\
 P(w,\vec e(u\mathbin{\dot{-}}w)).
 \tag{B}
\]
At zero substitute \(\vec e(u)\) for the free challenge in the
proved base formula \(P(0,\vec y)\). For the successor step, assume
\(Sw\leq u\); otherwise the implication is true. By (N) and the
recursion equation,
\[
 \vec e(u\mathbin{\dot{-}}w)
 =\vec b(w,\vec e(u\mathbin{\dot{-}}Sw)).
\]
Indeed write \(z=u\mathbin{\dot{-}}Sw\), so that
\(u\mathbin{\dot{-}}w=Sz\); the coefficient
\(u\mathbin{\dot{-}}Sz\) is \(w\) in this range, by the same
successor cancellation. The induction hypothesis (B), at its original
parameters, therefore supplies precisely the antecedent of the given
step theorem, with challenge
\(\vec e(u\mathbin{\dot{-}}Sw)\). That theorem yields the next
instance of (B). Set \(w=u\). Then \(u\mathbin{\dot{-}}u=0\)
and \(\vec e(0)=\vec y\), giving \(P(u,\vec y)\).
The only substitution at a changed challenge was into an already proved
step theorem, not into an unquantified induction assumption.
\end{proof}
\end{annexitem}


\begin{annexitem}{annex:p7.d.4}{thm:p07-dialectica-induction}{The full induction axiom}
\begin{proof}
Write the matrix of \(A(i)\) as \(P(i,\vec x,\vec y)\).
A positive witness for the antecedent of induction consists of
\(\vec a\) for \(A(0)\), and functionals \(\vec F,\vec G\)
for \(\forall i(A(i)\to A(Si))\). The latter matrix, at challenges
\(i,\vec x,\vec v\), is
\[
 P(i,\vec x,\vec G(i,\vec x,\vec v))\longrightarrow
 P(Si,\vec F(i,\vec x),\vec v).
 \tag{I}
\]
Define, by \cref{thm:p07-t-simultaneous-recursion},
\[
 \vec X(0)=\vec a,\qquad
 \vec X(Si)=\vec F(i,\vec X(i)).
\]
The positive witness for the conclusion is \(\lambda n.\vec X(n)\).
It has the required type and does not depend on the conclusion's
challenge \(\vec y\).

Given that challenge \((n,\vec y)\), define the backward sequence
\[
 \vec Y_n=\vec y,\qquad
 \vec Y_i=\vec G(i,\vec X(i),\vec Y_{i+1})\quad(i<n).
 \tag{Y}
\]
This is a System-T definition, not reverse recursion as a new primitive:
start at \(\vec y\), apply at step \(k\) the operation
\(\vec v\mapsto\vec G(n\mathbin{\dot{-}}Sk,
\vec X(n\mathbin{\dot{-}}Sk),\vec v)\), and take its value
after \(n\mathbin{\dot{-}}i\) steps. Tuple recursion supplies
all its components. The cutoff identities verified in the preceding
proof give exactly (Y) for \(i\leq n\). For indices above \(n\)
use the closed default tuple, so every displayed functional is total.

Put \(B_i=P(i,\vec X(i),\vec Y_i)\). These formulas have numerical
equality atoms, so \cref{thm:p07-t-matrix-decision} supplies their
Boolean tests. A bounded scan finds the first index \(i<n\) at
which \(B_i\land\neg B_{i+1}\), if one exists. Explicitly let
\(s(0)=0\), and at stage \(k+1\) retain \(s(k)\) if it is below
\(k\); otherwise take \(k\) when the new edge is bad, and \(k+1\)
when it is not. Thus \(s(k)=k\) is the sentinel for no bad edge.
This is numerical primitive recursion with the just constructed tests.
Open induction on \(k\), with the fixed parameters \(n,\vec y\),
proves
\[
\begin{split}
 s(k)<k&\ \Longrightarrow\ B_{s(k)}\land\neg B_{s(k)+1},\\
 s(k)=k&\ \Longrightarrow\ (B_0\to B_k).
\end{split}
\]
For the second implication the successor case with no old or new bad
edge uses \(B_k\to B_{k+1}\); the first implication either retains
its old witnessed edge or records the new one. These are precisely
the branches of the scan. Hence \(B_0\land\neg B_n\) implies
\(s(n)<n\) and a bad edge. Choose \(i=s(n)\) in that case and
zero otherwise.

Challenge the antecedent by \(\vec Y_0\) in its first conjunct and
by \((i,\vec X(i),\vec Y_{i+1})\) in its second. If the conclusion
matrix \(B_n\) is true there is nothing to prove. If it is false
but the first conjunct is true, then \(B_0\) is true and the scan
has selected a bad edge. By (Y) and the equation for \(\vec X\),
(I) at the selected challenge is exactly \(B_i\to B_{i+1}\),
which is false. Thus the two challenged antecedent conjuncts cannot
both hold unless \(B_n\) holds. When \(n=0\), the first conjunct
already is \(B_n\), and the unused second challenge may be arbitrary.
This is exactly the matrix of the induction axiom under the fixed
Dialectica clauses. All choices are typed conditional or recursion
terms, and every case verification is in the quantifier-free target.
\end{proof}
\end{annexitem}

\begin{annexitem}{annex:p7.d.5}{thm:p07-dialectica-soundness}{Soundness, rule by rule}
\begin{proof}
We verify the twelve groups of intuitionistic logical axioms and rules
in Section~2.1 of the held author preprint, using the handbook's
arithmetic axioms and induction schema. This fixes a concrete proof
presentation of HA for the extraction. Proofs are finite, substitutions
are capture avoiding, and numerical free parameters are retained
throughout. Write
\[
 A^{\mathrm D}=\exists\vec x\forall\vec y\,a(\vec x,\vec y),
 \quad B^{\mathrm D}=\exists\vec u\forall\vec v\,b(\vec u,\vec v),
 \quad C^{\mathrm D}=\exists\vec p\forall\vec q\,c(\vec p,\vec q).
\]
In particular a realizer \((F,G)\) for \(A\to B\) supplies the
proved target implication
\[
 a(\vec x,G(\vec x,\vec v))\longrightarrow
 b(F(\vec x),\vec v).
 \tag{R}
\]
All applications to tuples and all functions returning tuples stand for
componentwise typed applications and lists of functions. The following
checks are finite proof transformations, not semantic choices of
functionals.

\emph{1. Modus ponens.} Suppose a term \(\vec r\) realizes \(A\),
so the target proves \(a(\vec r,\vec y)\). Substitute
\(G(\vec r,\vec v)\) for \(\vec y\) and \(\vec r\) for
\(\vec x\) in (R). Modus ponens yields
\(b(F(\vec r),\vec v)\). The new positive witness is \(F(\vec r)\).

\emph{2. Composition.} Given \((F,G)\) for \(A\to B\) and
\((H,K)\) for \(B\to C\), take
\[
 F'(\vec x)=H(F(\vec x)),\qquad
 G'(\vec x,\vec q)=G(\vec x,K(F(\vec x),\vec q)).
\]
The first implication supplies the intermediate matrix
\[
 b(F(\vec x),K(F(\vec x),\vec q)).
\]
Applying the second implication yields the required \(c\)-matrix.

\emph{3. Contractions.} For \(A\lor A\to A\), a positive input is
\((d,\vec x_1,\vec x_2)\). Return
\(\operatorname{If}(d,\vec x_1,\vec x_2)\), and challenge both
input copies by the output challenge \(\vec y\). A true input
matrix forces \(d=0\) or \(d=1\) and the corresponding selected
matrix. For \(A\to A\land A\), return \((\vec x,\vec x)\)
and use the challenge selector at the end of
\cref{thm:p07-t-matrix-decision}. That lemma proves precisely the
required implication to the two challenged copies. The ordering of its
conditional branches is essential.

\emph{4. Injection and projection.} For \(A\to A\lor B\), return
\((0,\vec x,\vec d_B)\), where \(\vec d_B\) is the default
positive tuple for \(B\), and pass back the \(A\)-challenge.
For \(A\land B\to A\), return its \(A\)-witness and pair the
output challenge with a default challenge for \(B\). Propositional
introduction and elimination of conjunction verify the matrices.

\emph{5. Commutations.} For conjunction, exchange the positive tuples
and exchange the negative tuples. For disjunction, exchange those tuples
and replace its positive tag \(d\) by
\(\operatorname{If}(d,1,0)\). Whenever the input matrix is true,
its tag is zero or one, and the new tag selects exactly the same true
matrix in its exchanged position. Invalid tags make the input false.

\emph{6. Disjunction monotonicity.} From (R), interpret
\(C\lor A\to C\lor B\) by
\[
 (d,\vec p,\vec x)\longmapsto(d,\vec p,F(\vec x)).
\]
At output challenges \((\vec q,\vec v)\), challenge the input by
\((\vec q,G(\vec x,\vec v))\). The \(C\) branch is unchanged;
in the other branch use (R).

\emph{7. Currying and uncurrying.} A realizer for
\(A\to(B\to C)\) gives a matrix of the form
\[
 a(\vec x,H(\vec x,\vec u,\vec q))\to
 \bigl(b(\vec u,K(\vec x,\vec u,\vec q))\to
 c(F(\vec x,\vec u),\vec q)\bigr).
\]
The realizer for \(A\land B\to C\) is \(F\) with paired
negative output \((H,K)\). The two matrices are propositionally
equivalent. Conversely, separate that paired negative output, and
lambda-abstract the \(\vec u\) and \(\vec q\) arguments to
obtain the curried positive functions and the outer negative function.
All required types are exactly the arrow types specified by the
implication clause.

\emph{8. Falsity.} The matrix of \(\bottom\) is the same closed
false numerical equation and has no witness or challenge variables.
For \(\bottom\to C\), any default positive tuple suffices;
its matrix follows by propositional falsity elimination.

\emph{9. Universal introduction.} Suppose a realizer for
\(A\to B(z)\) gives \(F(z,\vec x)\) and
\(G(z,\vec x,\vec v)\), and \(z\) is not free in \(A\).
For \(A\to\forall zB(z)\), use positive tuple
\(\lambda z.F(z,\vec x)\) and negative tuple
\(G(z,\vec x,\vec v)\) at the new challenge \((z,\vec v)\).
The matrix is the original proved matrix after beta conversion.
The eigenvariable condition ensures that the input witness types and
matrix do not acquire an illicit dependence on \(z\).

\emph{10. Universal elimination.} For
\(\forall zA(z)\to A(t)\), a positive input is a tuple of
functions \(\vec f\). Return \(\vec f(t)\), and send the
challenge \(\vec y\) back as \((t,\vec y)\). Capture-free
substitution makes the two matrices identical.

\emph{11. Existential introduction.} For
\(A(t)\to\exists zA(z)\), return \((t,\vec x)\) and pass
back the unchanged challenge \(\vec y\).

\emph{12. Existential elimination.} From a realizer \((F,G)\) for
\(A(z)\to B\), where \(z\) is not free in \(B\), interpret
\(\exists zA(z)\to B\) by
\(F(z,\vec x)\) on positive input \((z,\vec x)\), and return
\(G(z,\vec x,\vec v)\) as its negative challenge. Again the
matrix is the given theorem. The eigenvariable condition is precisely
what makes the output a witness for the unchanged \(B\).

\emph{Equality.} Reflexivity is its own matrix. Atomic replacement
\(s=t\to(\chi[s/z]\to\chi[t/z])\), for numerical terms, is
its own quantifier-free matrix and follows from target equality
replacement. These are the equality axioms of the logical presentation;
replacement in compound formulas is obtained by their ordinary logical
derivations. Thus it does not introduce a new higher-type equality
principle into the interpretation.

\emph{Arithmetic.} Define
\[
 \operatorname{add}(a,b)=R_0\,a\,(\lambda k,r.Sr)\,b,
 \qquad
 \operatorname{mul}(a,b)=R_0\,0\,
                  (\lambda k,r.\operatorname{add}(r,a))\,b.
\]
Their equations are exactly the four addition and multiplication axioms
of the selected arithmetic language. Successor injectivity and its
nonzero value are target axioms. The remaining Robinson axiom is not
quantifier free: \(n\ne0\to\exists m\,n=Sm\). Its positive
witness is \(p(n)\), with predecessor as previously defined. Open
induction proves \(n=0\lor n=Sp(n)\), so the interpreted implication
is a target theorem. Universal closures use the preceding quantifier
constructions. Finally every induction axiom is realized by
\cref{thm:p07-dialectica-induction}. This covers the handbook's full
arithmetic axiom list, not only its quantifier-free part.

Induction on the given finite HA derivation assembles the indicated
terms and target proofs. Intermediate premises can have numerical
parameters absent from their conclusion. Retain them while performing
the specified eigenvariable abstractions and substitutions; then replace
any remaining extraneous numerical parameter by zero in the terms and
target proof. The conclusion formula is independent of those parameters,
so theorem substitution preserves its interpreted matrix. Negative
challenge variables occur only in the matrix, not freely in its positive
tuple. For a closed conclusion the resulting positive tuple is closed. Every
construction is a typed term substitution, a finite logical proof
construction, or one of the specified recursion constructions. This
proves the theorem and its effectivity assertion.
\end{proof}
\end{annexitem}

\begin{annexitem}{annex:p7.d.6}{thm:p07-dialectica-consistency}{The syntactic consistency implication}
\begin{proof}
Use explicit finite proof codes for the source presentation and for the
target just specified. Checking a displayed typed term, an axiom
instance, a single beta or recursor equation, a substitution, or an
inference is primitive recursive. An induction instance is checked by
its finite formula and displayed substitutions, not by deciding its
truth. The preceding proof gives a primitive-recursive transformation
of a source derivation into a target derivation: each inference replaces
already constructed terms and proof fragments in fixed finite templates.
There is no appeal to normalization, proof search, or an unbounded
semantic test. The recursion is on the finite source proof code.

Primitive recursive arithmetic formalizes the verification of this
transformation by quantifier-free induction on that code. In particular,
for the proof predicates of these fixed presentations it proves
\[
 \operatorname{Prf}_{\mathsf{HA}}(p,\ulcorner\bottom\urcorner)
 \ \longrightarrow\
 \operatorname{Prf}_{\mathsf T_{\mathrm{qf}}}
       (f(p),\ulcorner\bottom\urcorner).
\]
Here \(f\) is the specified syntactic transformer. The interpretation
of \(\bottom\) has no positive or negative variables and is exactly
\(\bottom\), so no additional truth-reflection principle is used.
Taking the contrapositive of the existential assertion about proof
codes yields
\(\mathsf{Con}(\mathsf T_{\mathrm{qf}})\to
\mathsf{Con}(\mathsf{HA})\).
This is a relative syntactic consistency statement in the stated
arithmetical base, not a proof of either consistency assertion in HA.
\end{proof}
\end{annexitem}
